\chapter{Aggregation theory}\label{ch:aggregation}

Part IV asks a question the earlier parts could not: at what grain does this bank's data carry information? A grain is
the unit one row of a fact describes, a customer in a month, a market on a day, an agent in an hour. Aggregating to a
coarser grain removes noise and also removes detail; it can create correlations that hold for no individual and hide
signals that only the individual carries. This chapter sets out the rules the granularity series built its aggregates
by, the statistical profile of each grain, and the traps of aggregation, the ecological fallacy and exposure, before
the pooled ledger of every target the series evaluated. The chapters that follow examine customers
(\cref{ch:customer-dynamics}), operations (\cref{ch:operations}) and products and campaigns
(\cref{ch:products-and-campaigns}) grain by grain.

\section{Grains, stars and the tests of an aggregate}\label{sec:ag2-star}

The series designs its aggregates as a dimensional star: a bus matrix of business processes against conformed dimensions,
and one fact table per grain \autocite{kimball2013}. Every aggregate must pass the same tests. Its \emph{grain is unique}: no rows share the key that defines the grain. It \emph{reconciles}: summed over its rows, every additive measure equals the
same measure counted on the atomic rows. And it is \emph{dense} where it must be: every unit appears in every period, with
an explicit zero when nothing happened, because a rolling average over a sparse table silently skips the quiet periods. The
stored build checks pass all of them: $\est{ch14.checks-day-ok}{value}$ of $\est{ch14.checks-day}{value}$ checks on the
$\est{ch14.models-day}{value}$ models of the day and month star, and $\est{ch14.checks-hour-ok}{value}$ of
$\est{ch14.checks-hour}{value}$ on the $\est{ch14.models-hour}{value}$ models of the hour star.

Measures are stored so that they add up. A rate is stored as its numerator and denominator, never as a ratio, because the
rate of a coarser grain is a ratio of sums,
\begin{equation}\label{eq:ag2-ratio}
  R = \frac{\sum_i n_i}{\sum_i d_i} = \sum_i \frac{d_i}{\sum_j d_j}\, r_i ,
\end{equation}
a weighted mean of the finer rates $r_i = n_i / d_i$ with weights proportional to their denominators, and not their
plain average. The plain average gives a quiet day the same weight as a busy one and changes when the same data are cut
differently, which is the arithmetic behind Simpson's paradox.

\section{Zeros and dispersion}\label{sec:ag2-profile}

Properties of a grain's main count decide which models suit it (\cref{tab:ch14-grains}). The share of zeros says
whether most units do nothing in most periods; the ratio of variance to mean says whether units differ beyond chance,
since a Poisson count has a ratio of one and heterogeneity between units raises it \autocite{cameron2013}.

\input{../build/generated/tables/ch14-grains}

Zeros alone do not call for a zero-inflated model \autocite{lambert1992}; zeros in excess of what the mean implies do. A
Poisson count with mean $\mu$ is zero with probability $e^{-\mu}$. Customer-month contacts have a mean of
$\est{ch14.cm-contact-mean}{value}$, so $\est{ch14.cm-contact-pois-zero}{value}\,\%$ zeros are expected and
$\est{ch14.cm-contact-zero}{value}\,\%$ are observed, with a variance-to-mean ratio of $\est{ch14.cm-contact-vm}{value}$;
complaints expect $\est{ch14.cm-complaint-pois-zero}{value}\,\%$ and show $\est{ch14.cm-complaint-zero}{value}\,\%$, ratio
$\est{ch14.cm-complaint-vm}{value}$. Both are Poisson to the precision shown: every customer contacts the bank and complains at
the same rate, so there is nothing at this grain that distinguishes one customer from another. Transactions differ. With a
mean of $\est{ch14.cm-tx-mean}{value}$ a Poisson count would be zero in $\est{ch14.cm-tx-pois-zero}{value}\,\%$ of
customer-months; the data show $\est{ch14.cm-tx-zero}{value}\,\%$ and a ratio of $\est{ch14.cm-tx-vm}{value}$. Customers do
differ in how much they transact, and \cref{sec:ag2-exposure} shows why. Product-months are idle
$\est{ch14.pm-tx-zero}{value}\,\%$ of the time. At the market-day grain there are no zeros
($\est{ch14.md-tx-zero}{value}\,\%$), and the variance-to-mean ratio of $\est{ch14.md-tx-vm}{value}$ measures differences
between markets and weekdays, not noise; once those are removed, the daily residuals are independent
(\cref{sec:cp-residual}).

\paragraph{Hour grains.} Below the day the question is whether a grain is populated at all. If events in a cell arrive as
a Poisson process with mean $\lambda$, the cell is occupied with probability $1 - e^{-\lambda}$. For every hour grain the
series considered except one, the observed occupancy matches that prediction within
$\est{ch14.occ-gap-other}{value}$ percentage points (\cref{tab:ch14-sparsity}): the events of a customer, a product, a branch
or an agent within the hour fall at random, with no clustering a model could use. The exception, market by channel by
hour, is occupied $\est{ch14.occ-gap-max}{value}$ points less than Poisson predicts, because some channels are closed at
some hours: structural zeros, the one kind that a zero-inflated model is for. Only $\est{ch14.viable-hour}{value}$ of the
$\est{ch14.hour-grains}{value}$ hour grains carry enough events per cell to be worth storing; a customer is active in
$\est{ch14.customer-hour-occ}{value}\,\%$ of their hours.

\input{../build/generated/tables/ch14-sparsity}

\section{The ecological fallacy}\label{sec:ag2-ecological}

A correlation between aggregates says nothing, by itself, about the units inside them \autocite{robinson1950}. The law of
total covariance makes the reason precise. For measures $X$ and $Y$ of units grouped into cells $G$,
\begin{equation}\label{eq:ag2-covariance}
  \operatorname{Cov}(X, Y) = \mathbb{E}\bigl[\operatorname{Cov}(X, Y \mid G)\bigr]
  + \operatorname{Cov}\bigl(\mathbb{E}[X \mid G],\ \mathbb{E}[Y \mid G]\bigr),
\end{equation}
and an analysis of cell means sees only the second term, which a driver common to $X$ and $Y$ at the cell level can make
large while the first is zero. The series measured this with one pair of measures, digital errors and contacts. Across
market-days the rank correlation is $\est{ch14.eco-md-raw}{value}$; once the weekday is removed it is
$\est{ch14.eco-md-adj}{value}$. Across customer-months it is $\est{ch14.eco-cm}{value}$, across customers' lifetimes
$\est{ch14.eco-life}{value}$, and $\est{ch14.eco-life-adj}{value}$ after controlling for how long each customer was active.
The aggregate correlation was made entirely by the calendar: busy days are busy for every process. Customers with errors do
not call more. Every co-movement shown on a dashboard of aggregates is a hypothesis about units, to be tested at the unit
grain after the common drivers are removed.

\section{Exposure}\label{sec:ag2-exposure}

Dormancy, a month without transactions, is the one customer target the series found learnable, and its signal has a simple
source. Suppose each product a customer holds is used in a month independently with probability $p$. Then a customer
holding $k$ products is dormant with probability
\begin{equation}\label{eq:ag2-exposure}
  \Pr(\text{dormant} \mid k) = c\,(1 - p)^{k},
\end{equation}
with $c$ absorbing activity not tied to a held product. On a log scale this is a straight line in $k$. Fitted by weighted
least squares to the dormant share of customer-months by products held, the line has ratio $r = \est{ch14.exp-r}{value}$,
intercept $c = \est{ch14.exp-c}{value}$ and $R^2 = \est{ch14.exp-r2}{value}$, deviating from the data by at most
$\est{ch14.exp-gap}{value}$; the products are used in $\est{ch14.exp-p}{value}$ of product-months, so independence predicts a
ratio of $1 - p$, which matches. The dormant share falls from $\est{ch14.exp-dormant-1}{value}$ for one product to $\est{ch14.exp-dormant-8}{value}$ for the most products held, entirely by arithmetic. A customer-level dormancy model therefore learns, above
all, how many products a customer holds: exposure, not behaviour. This is also the heterogeneity behind the excess zeros of
\cref{sec:ag2-profile}, and it is the yardstick any claim of behavioural signal must beat.

\section{The pooled target ledger}\label{sec:ag2-ledger}

The series evaluated $\est{ch14.ledger-targets}{value}$ targets across $\est{ch14.ledger-grains}{value}$ grains:
classifications judged by their AUC, survival models by their concordance \autocite{harrell1982}, forecasts by the mean
absolute scaled error against the seasonal-naive forecast \autocite{hyndman2006}. Each carries a one-sided $p$-value: for an AUC, against chance from the Hanley--McNeil standard error \autocite{hanley1982}; for a forecast, the Diebold--Mariano test
of equal absolute loss against the naive forecast \autocite{diebold1995}, with the small-sample correction of
\textcite{harvey1997}. The ledger pools all of them and controls the false-discovery rate (\cref{eq:bh}); recomputed here, its
$q$-values agree exactly with the stored ones ($\est{ch14.ledger-q-reproduced}{value}$ for agreement). Of the
$\est{ch14.ledger-tested}{value}$ targets, $\est{ch14.ledger-discoveries}{value}$ survive at a false-discovery rate of $\est{ch14.fdr-level}{value}\,\%$ (\cref{tab:ch14-ledger,tab:ch14-discoveries}).

\input{../build/generated/tables/ch14-ledger}

\input{../build/generated/tables/ch14-discoveries}

The discoveries are of a few kinds: the dormancy and lapse targets that exposure explains, a few daily forecasts that beat the
weekly naive forecast (taken up in \cref{ch:operations}), and campaign differences (\cref{ch:products-and-campaigns}). One
deserves a warning: the early-warning rule that customers declined for insufficient funds go dormant has a relative risk below one, so it warns in the wrong direction, and the series itself judged the association immaterial. The direction has a mechanical reason: a customer must transact to be declined, so declined customers are active by construction, another face of exposure. And one caution applies to the forecasts: the Diebold--Mariano variance is computed
here without any autocorrelation term, which is right for one-step forecast errors that are uncorrelated and optimistic
otherwise.

\section{Verdict and implications}\label{sec:ag2-verdict}

The aggregates are sound by construction: unique grains, reconciled measures, dense grids, rates as ratios of sums. What
they reveal about this bank is mostly the absence of structure. Contacts and complaints are Poisson at the customer-month
grain, events fall at random within hours, an aggregate correlation between errors and contacts is the calendar's, and the
learnable customer signal is the arithmetic of exposure. A small set of grains carries information worth promoting, and
the chapters of this part name them. Rules from this chapter apply to every later use of the aggregates: test a
co-movement at the unit grain before believing it, and compare any behavioural signal with the exposure it might merely
reflect.

\section{Limitations}\label{sec:ag2-limits}

The zero and dispersion profile is recomputed here on processing months and days, a close stand-in for the series' local
calendar; the build checks, the hour-grain sparsity, the correlations, the exposure table and the ledger are the series'
stored results. The ledger's $p$-values come from the series' own models, which this chapter does not refit; their
interpretation depends on the caveats noted for each test.
